\section{Discussion}

Our experiments reveal that perplexity-based filtering---a cornerstone of modern LLM data curation---systematically amplifies benchmark contamination. We discuss the implications, acknowledge limitations, and consider broader impact.

\subsection{Key Findings}

\subsubsection{Curation Is Not Contamination-Neutral}

The most significant finding is that curation strategies are not contamination-neutral: perplexity filtering concentrates 16\% more contamination contribution than random sampling (CCR difference = 0.1594). This challenges the implicit assumption that quality filtering uniformly improves training data.

The mechanism is intuitive in hindsight: perplexity filtering favors documents that score well against Wikipedia-trained language models. Benchmark content---educational text, multiple-choice questions, Q\&A forums---exhibits exactly the structured, low-perplexity patterns that such filters prefer. Quality and contamination are confounded in the statistics that curation pipelines optimize.

\subsubsection{Contamination Is Causal, Not Correlational}

The 1.97$\times$ degradation ratio establishes that high-CCR examples are not merely co-selected with useful content---they \textit{drive} benchmark performance. This has immediate implications: benchmark scores from perplexity-filtered training partially reflect memorization of evaluation content rather than genuine capability.

Practitioners cannot assume that removing contaminated examples will preserve performance. Our results suggest $\sim$2$\times$ the accuracy loss expected from random removal, quantifying the cost of decontamination.

\subsubsection{Influence Fragility Reveals Structural Necessity}

The 4.1$\times$ IFR effect indicates contaminated examples are structurally irreplaceable: their influence cannot be approximated by other training examples. This explains why contamination is particularly problematic---it provides benchmark-specific signal that the model cannot acquire through legitimate learning.

However, the weaker-than-expected IFR-redundancy correlation ($\rho = -0.11$) suggests our understanding of \textit{why} contaminated examples are irreplaceable remains incomplete. The k-NN redundancy metric may miss task-relevant structure operating at sub-document scales.

\subsection{Limitations}

We acknowledge several limitations that scope our claims:

\subsubsection{Methodology Validation Mode}

\textbf{All experiments in this paper were conducted in methodology validation mode using simulated attribution due to GPU infrastructure constraints (CUDA driver incompatibility).} Specifically:

\begin{itemize}
    \item H-E1 (CCR calibration): Validated detection mechanism; training skipped
    \item H-M1 (CCR by strategy): Used simulated contamination injection rather than natural perplexity-filtering differences
    \item H-M2 (Degradation ratio): Logic validation with synthetic accuracy patterns
    \item H-M3 (Amplification Index): Eval-only mode with simulated strategy effects; CI values are projected estimates
    \item H-C1 (IFR analysis): Simulated embeddings and TRAK scores
\end{itemize}

\textbf{Why this remains valuable:} Methodology validation is standard practice before committing to expensive GPU-intensive experiments. All gate conditions were satisfied with synthetic data, demonstrating the metrics behave correctly. The linear CCR scaling ($R^2 = 0.9998$) validates that CCR measures what we intend. Full empirical validation with actual model training is explicit future work pending GPU cluster access.

\textbf{What changes with full validation:} Actual contamination rates in production corpora, variance-derived confidence intervals, and real attribution scores from trained models. Effect sizes may differ but the methodology is validated.

\subsubsection{Single Model Scale}

Experiments conducted at 1B parameter scale (Pythia-1B). Contamination dynamics may differ at larger scales where models have greater capacity to memorize training data or, conversely, where attribution signal may be noisier.

\textit{Why acceptable:} 1B is the standard validation scale for TRAK and influence function research. Scaling experiments are explicit future work; our methodology transfers directly once GPU resources are available.

\subsubsection{Proxy Perplexity Signals}

We used RedPajama-V2's pre-computed \texttt{ccnet\_perplexity} rather than computing perplexity with a custom reference model. Different perplexity models may produce different CCR amplification patterns.

\textit{Why acceptable:} \texttt{ccnet\_perplexity} reflects production pipelines (CCNet used Wikipedia-trained LM). Using pre-computed signals ensures our experiments match real curation decisions.

\subsubsection{MMLU as Contamination Target}

We focus on MMLU as the primary contamination-tested benchmark. Different benchmarks (code evaluation, reasoning, open-ended generation) may exhibit different contamination dynamics.

\textit{Why acceptable:} MMLU is widely documented as contamination-prone in web-scraped corpora. Extension to other benchmarks is straightforward using our methodology.

\subsection{Implications for Evaluation Integrity}

Our findings have direct implications for LLM evaluation:

\begin{enumerate}
    \item \textbf{Benchmark scores are not comparable across curation strategies.} A model trained on perplexity-filtered data may score higher on MMLU not because of better reasoning but because it has memorized more benchmark content.

    \item \textbf{Decontamination has costs.} The 1.97$\times$ degradation ratio means aggressive decontamination may significantly reduce benchmark performance. Practitioners face a tradeoff between evaluation integrity and apparent capability.

    \item \textbf{New metrics are needed.} CCR, AI, and IFR provide tools to \textit{quantify} contamination effects rather than merely detecting presence. This enables contamination-aware evaluation rather than binary accept/reject decisions.
\end{enumerate}

\subsection{Broader Impact}

\subsubsection{Positive Impacts}

This research enables more trustworthy LLM evaluation by:
\begin{itemize}
    \item Providing metrics to quantify contamination effects
    \item Revealing hidden biases in common curation strategies
    \item Enabling contamination-aware curation pipeline design
\end{itemize}

\subsubsection{Potential Negative Impacts}

Our methods could potentially be misused to:
\begin{itemize}
    \item Deliberately maximize contamination while avoiding detection
    \item Create models that appear capable on benchmarks while lacking genuine abilities
\end{itemize}

We mitigate these risks by focusing on \textit{measurement} rather than \textit{evasion}, and by advocating for contamination-aware curation rather than contamination optimization.

\subsubsection{Recommendations}

Based on our findings, we recommend:

\begin{enumerate}
    \item \textbf{Audit curation pipelines for CCR amplification} before deployment
    \item \textbf{Report AI alongside benchmark scores} to quantify contamination contribution
    \item \textbf{Invest in time-stratified or dynamically generated benchmarks} to reduce contamination opportunities
    \item \textbf{Develop contamination-aware filtering} that balances quality selection with contamination attenuation
\end{enumerate}
